\paragraph{The descended class and its floor restriction.}
In this paragraph and the next three, write \(X=X_i\), \(Z=Z_i\), and
\(\alpha=\alpha_i\).  We use
\(dd^c=\frac{i}{2\pi}\partial\bar\partial\), the normalization for
which \(dd^c\phi\) is the normalized Chern curvature of a metric
locally written as \(e^{-\phi}\).  A fixed positive rescaling of potentials gives the
\(i\partial\bar\partial\) convention in the cited analytic inequalities.

Both ends of \(c\) are normal compact spaces in class \(\mathcal C\)
with rational singularities, and \(\alpha\) is zero on every contracted
curve.  The Bott--Chern descent
\cite[Lemma~2.11]{DasHacon2024} therefore gives a class \(\beta\) on
\(Z\) and a smooth real closed representative \(\beta_0\), with local
smooth potentials, such that
\begin{equation}\label{eq:outer-descended-class}
 c^*\beta=\alpha\quad\text{in }H^{1,1}_{BC}(X).
\end{equation}
Here and below the group is the local-potential Bott--Chern group.
Its smooth representative description, including adjustment by a global
smooth \(dd^c\)-potential, is
\cite[Definition~3.1 and Remark~3.2]{HoeringPeternell2016}.

Let \(i_W:W\hookrightarrow Z\) be the embedded floor image and put
\(\delta=i_W^*\beta\).  The actual identity
\(c|_T=i_W\circ\varphi\), \eqref{eq:floor-contraction-class}, and
\eqref{eq:outer-descended-class} give
\begin{equation}\label{eq:outer-floor-image-class}
 \varphi^*\delta=\alpha|_T=\varphi^*[\omega_W]
 \quad\text{in }H^{1,1}_{BC}(T).
\end{equation}
Thus the pullback of \(\delta-[\omega_W]\) is zero and hence nef.
Apply \cite[Lemma~2.38]{DasHaconPaun2024} only to
\(\varphi:T\to W\): it is proper and surjective, and both spaces are
normal compact K\"ahler.  The lemma makes \(\delta-[\omega_W]\) nef.
For a smooth representative \(e\) of that difference, choose a nef
approximation \(e+dd^c u\geq-\omega_W/2\).  Then
\(\omega_W+e+dd^c u\) is a K\"ahler representative of \(\delta\).
In particular its restriction gives a positive current in the class
\(\beta|_V\) and a positive top integral for every positive-dimensional
irreducible \(V\subset W\).  Only the nefness of
\(\delta-[\omega_W]\) enters this deduction.

\paragraph{Nefness on the target.}
For every positive-dimensional irreducible reduced compact subspace
\(V\subset Z\), we first produce a current
\begin{equation}\label{eq:outer-restriction-current}
 T_V=\beta_0|_V+dd^c q_V\geq0
\end{equation}
with a global real distribution \(q_V\) on \(V\).  Currents and their
positivity on the pure-dimensional reduced \(V\) are defined by duality
with smooth test forms from local ambient embeddings; see
\cite[Section~1, Definitions~1.1--1.2, pp.~14--15]{Demailly1985}.
We will use exactly this positive-current meaning of pseudoeffectivity
on a possibly nonnormal subspace.  The preceding K\"ahler representative
of \(\delta\) gives \eqref{eq:outer-restriction-current} for
\(V\subset W\), with a smooth potential.

Suppose \(V\not\subset W\), and let \(V'\subset X\) be the closure
of the inverse image of \(V\setminus W\).  The isomorphism
\(X\setminus T\simeq Z\setminus W\) makes \(V'\) irreducible and
its map to \(V\) proper and bimeromorphic.  Choose a projective
resolution \(r:U\to V'\) with smooth source.  The restricted K\"ahler
form makes \(V'\) a compact K\"ahler space.  A relatively positive
metric for the projective resolution, plus a sufficiently large pullback
of this form, makes \(U\) a compact K\"ahler manifold.  This metric
construction applies to the possibly nonnormal \(V'\).

Let \(\nu:V^\nu\to V\) be normalization.  The dominant map from the
normal \(U\) factors through \(\nu\), giving a proper bimeromorphic
map \(p:U\to V^\nu\), hence a modification between normal compact
spaces.  For the other map \(j:U\to X\), functoriality of the actual
maps and \eqref{eq:outer-descended-class} gives
\[
 p^*\nu^*(\beta|_V)=j^*\alpha.
\]
The class on the right is nef.  Indeed, pull back the smooth nef
approximations from \(X\); on compact \(U\), the pullback of their
reference form is bounded above by a fixed multiple of a K\"ahler form
\(\eta\) on \(U\), so rescaling the error gives the nef inequalities.

Set \(\gamma=p^*\nu^*(\beta_0|_V)\).  Choose smooth approximations
\(\gamma+dd^c f_n\geq-n^{-1}\eta\).  The closed positive currents
\(\gamma+dd^c f_n+n^{-1}\eta\) have bounded \(\eta\)-mass by
Stokes' theorem.  Weak compactness gives a positive closed limit in
\([\gamma]\).  The \(\partial\bar\partial\)-lemma on the compact
K\"ahler manifold \(U\), and the global potential description of a
positive current, give a global quasi-plurisubharmonic \(q_U\) with
\[
 \gamma+dd^c q_U\geq0.
\]
These are also the closed-cone statement following Definition~1.6 and
the potential description in (3.1) and its following paragraph in
\cite[pp.~1253 and 1260]{DemaillyPaun2004}.  Apply
\cite[Corollary~2.32 and its proof, p.~18]{DasHaconPaun2024} to \(p\), with target form
\(\nu^*(\beta_0|_V)\) and constant zero in its current inequality.
The locally integrable, locally bounded above \(q_U\) supplies the
required potential.  We obtain a global quasi-plurisubharmonic
\(q_\nu\) satisfying
\begin{equation}\label{eq:outer-normalized-current}
 T_\nu=\nu^*(\beta_0|_V)+dd^c q_\nu\geq0.
\end{equation}
This applies the normal-modification descent to the displayed
local-potential inequality on the still general normal compact target.

It remains to push through the finite normalization.  Locally write
\(\beta_0|_V=dd^c h\) with \(h\) smooth.  The function
\(\nu^*h+q_\nu\) is locally integrable and locally bounded above, and
its Hessian is \(T_\nu\geq0\).  On normal \(V^\nu\), its upper
regularization is plurisubharmonic.  The finite trace is weakly
plurisubharmonic and its Hessian is the positive current \(\nu_*T_\nu\)
by \cite[Corollary~1.11 and Proposition~1.13(a), p.~22]{Demailly1985}.
Since \(\nu\) has generic degree one, change of variables off the
proper analytic exceptional sets gives
\(\nu_*\nu^*(\beta_0|_V)=\beta_0|_V\) as currents; those sets have
measure zero for the smooth top-degree integrands.  Proper push of
distributions commutes with \(dd^c\).  Hence
\begin{equation}\label{eq:outer-finite-trace}
 T_V=\nu_*T_\nu
     =\beta_0|_V+dd^c(\nu_*q_\nu)\geq0.
\end{equation}
The traced distributions are locally integrable and glue because
\(q_\nu\) is global.  On a locally reducible \(V\), these local
potentials may be only weakly plurisubharmonic; the argument below uses
their associated positive currents.

We now apply the sufficient direction of the restriction criterion
\cite[Theorem~2.36 and Remark~2.37]{DasHaconPaun2024}, spelling out its
singular-current input.  Proposition~3.3(iv), p.~1262, and the following
paragraph, pp.~1262--1263, of \cite{DemaillyPaun2004}
applies to a compact complex space and a smooth class containing a
closed positive current.  It gives nefness if the restrictions to the
irreducible components of every positive Lelong level set are nef.
Induct on the dimension of each irreducible \(V\subset Z\); points are
automatic.  For a positive-dimensional \(V\), use
\eqref{eq:outer-restriction-current}.  The Siu analyticity assertion
accompanying that proposition makes each positive Lelong level set
analytic, and it is proper.  Indeed, if a fixed positive level filled \(V\), pack
disjoint radius-\(r\) balls in a coordinate ball of \(V_{\rm reg}\).
The Lelong mass lower bound on each ball would force the locally finite
mass of the \((1,1)\)-current to grow at least as a positive multiple
of \(r^{-2}\) when \(r\) tends to zero.  This is impossible.  The level
set components therefore have smaller dimension, and their restricted
classes are nef by induction.  The proposition makes \(\beta|_V\) nef.
Taking \(V=Z\) proves analytic nefness of \(\beta\).  This uses the
positive-current formulation above also on nonnormal subspaces, where
the local potentials may be only weakly plurisubharmonic.

\paragraph{Bigness from the supported section.}
Use the effective actual representative already present in
\eqref{eq:running-supported-representative}.  Choose \(m>0\) clearing
the indices and an isomorphism of holomorphic lines
\[
 L=\OO_X(mA_i)\simeq\OO_X(mP_i).
\]
The canonical section of the right side gives a nonzero holomorphic
section \(s\) of \(L\) with divisor \(mP_i\).  Choose a smooth metric
\(h\) on \(L\).  In a local frame write \(h=e^{-\phi}\), \(s=f\),
and put
\[
 \vartheta=\frac1m dd^c\phi,\qquad
 \ell=\frac1m\log|s|_h^2.
\]
The first formula defines the normalized global curvature form.
The second is a global quasi-plurisubharmonic function, locally
\(m^{-1}(\log|f|^2-\phi)\).  The holomorphic \(f\) extends in a local
ambient embedding and is not identically zero on any nonempty open
subset of that chart in \(X\).  Local integrability is Proposition~1.8
of \cite[p.~19]{Demailly1985}, and
the unnumbered prose after its proof gives the current positivity
\[
 \vartheta+dd^c\ell=\frac1m dd^c\log|f|^2\geq0.
\]
The smooth form \(\vartheta+\omega_i\) represents
\(\alpha=c^*\beta\), so the smooth local-potential representative
description supplies a global smooth \(v\) with
\(\vartheta+\omega_i=c^*\beta_0+dd^c v\).  Thus
\begin{equation}\label{eq:outer-supported-current}
 c^*\beta_0+dd^c(v+\ell)\geq\omega_i.
\end{equation}
This is a global quasi-plurisubharmonic potential obtained from the
supported section at this running step.

Fix any smooth positive Hermitian form \(g_Z\) on \(Z\) in local
ambient embeddings.  Local holomorphic coordinate lifts of \(c\) make
\(c^*g_Z\) semipositive in those embeddings.  A finite relatively
compact cover of \(X\) and comparison with the positive definite
ambient representatives of \(\omega_i\) give one \(C>0\) with
\(c^*g_Z\leq C\omega_i\).  Corollary~2.32 of
\cite{DasHaconPaun2024} applies to the normal modification \(c\) and
\eqref{eq:outer-supported-current}.  It gives a global
quasi-plurisubharmonic, hence \(L^1\), function \(v_Z\) with
\begin{equation}\label{eq:outer-current-big}
 \beta_0+dd^c v_Z\geq C^{-1}g_Z.
\end{equation}
This is the required dominating current in the selected target class.

\paragraph{Top intersections and the other floor.}
Let \(V\subset Z\) be irreducible of dimension \(k>0\).  The case
\(V\subset W\) was proved using the K\"ahler class \(\delta\).
Otherwise use its strict transform \(V'\) above.  The proper integration
cycle identity \(c_*[V']=[V]\), projection formula, and
\eqref{eq:outer-descended-class} give
\begin{equation}\label{eq:outer-strict-transform}
 \int_{V_{\rm reg}}\beta_0^k
   =\int_{V'_{\rm reg}}\alpha^k.
\end{equation}
The right side denotes the integral of any smooth representative of
\(\alpha\); Stokes' theorem makes it independent of that choice.
Neither subspace needs to be normal.

If \(V'\not\subset\Supp P_i\), take the projective resolution
\(r:U\to V'\) with smooth K\"ahler source used above.  The pullback of
\(s\) is nonzero and defines
the effective rational divisor \(E=m^{-1}\operatorname{div}(r^*s)\)
on \(U\).  Put \(a=r^*(\alpha|_{V'})\) and
\(w=r^*[\omega_i|_{V'}]\).  The class \(a\) is nef, \(w\) has a
semipositive smooth representative, and the actual section gives
\(a-w=c_1(E)\).  Therefore
\begin{equation}\label{eq:outer-outside-support}
 \int_U a^k-\int_U w^k
 =\sum_{j=0}^{k-1}c_1(E)\cdot a^{k-1-j}w^j\geq0.
\end{equation}
For each summand, choose a smooth representative of
\(a+\varepsilon[\eta]\) which is semipositive, where \(\eta\) is a
fixed K\"ahler form on \(U\).  Restrict it and the semipositive
representative of \(w\) to every effective divisor component, integrate,
and let \(\varepsilon\) tend to zero.
This proves its nonnegativity, including the degree statement when
\(k=1\).  Projection formula gives
\(\int_U w^k=\int_{V'_{\rm reg}}\omega_i^k>0\).  Equations
\eqref{eq:outer-strict-transform}--\eqref{eq:outer-outside-support}
give the desired strict positivity.

It remains that \(V'\) lies in a component \(T'\) of \(\Supp P_i\).
It meets \(X\setminus T\), so \(T'\neq T\).  The actual floor
specialization in Subsection~\ref{subsec:actual-floor-split} gives
\begin{equation}\label{eq:outer-other-floor}
 g:T'\to\overline W,\qquad \alpha|_{T'}=g^*\kappa,
\end{equation}
where \(g\) is projective with connected fibers, its source and target
are normal compact K\"ahler, and \(\kappa\) is a K\"ahler class.
The equality is in the local-potential Bott--Chern group.
Every irreducible curve in a \(g\)-fiber has \(\alpha\)-degree zero,
so \eqref{eq:global-support-zero-ray} puts its class in \(R_i\), and
\(c\) contracts it.

It follows that \(h=c|_{T'}\) is pointwise constant on every full
\(g\)-fiber.  To see this also for nonreduced fibers, their reductions
are connected projective varieties, possibly reducible.  If the image
of such a reduction were
not a point, one irreducible component \(Q\) would have nonconstant
image; otherwise its image would be a connected finite set.  At a smooth
point of \(Q\) where a local coordinate of \(h\) has nonzero
differential, general hyperplanes through that point cut an irreducible
projective curve component whose tangent is not in the differential's
kernel.
That curve has nonconstant image, a contradiction.  If \(\dim Q=1\),
take \(Q\) itself.

There is consequently a holomorphic factorization
\begin{equation}\label{eq:outer-factorization}
 h=b\circ g,\qquad b:\overline W\to Z.
\end{equation}
For clarity, the specialization gives
\(g_*\OO_{T'}=\OO_{\overline W}\); this also follows from analytic
Stein factorization, connected fibers, and the normality of the two
spaces.  Proper surjectivity makes \(g\) a closed quotient map, so the
pointwise constancy first defines a continuous \(b\).  For an open
\(O\subset Z\) embedded in a polydisc, restrict to the open
\(b^{-1}(O)\).  The coordinate functions of \(h\) on its full inverse
image descend through \(g_*\OO_{T'}=\OO_{\overline W}\) to a
holomorphic tuple.  It lies pointwise in the embedded \(O\);
its defining ideal therefore vanishes on the reduced open \(b^{-1}(O)\).
The resulting local holomorphic maps glue to
\eqref{eq:outer-factorization}.  This factorization uses pointwise
constancy on the underlying full fibers.

Put \(H=g(V')\), a closed irreducible analytic subspace of
\(\overline W\).  Since \(c|_{V'}\) is bimeromorphic onto \(V\),
\eqref{eq:outer-factorization} yields
\[
 k=\dim c(V')=\dim b(H)\leq\dim H\leq\dim V'=k.
\]
Thus \(g|_{V'}\) is generically finite onto \(H\), of some integer
degree \(d>0\).  The projection formula using
\eqref{eq:outer-other-floor} now gives
\begin{equation}\label{eq:outer-positive-intersections}
 \int_{V_{\rm reg}}\beta_0^k
 =\int_{V'_{\rm reg}}(g^*\kappa)^k
 =d\int_{H_{\rm reg}}\kappa^k>0.
\end{equation}
The last integral is a positive K\"ahler volume.  The one-way
factorization \(h=b\circ g\) and this degree calculation give the needed
implication in \cite[proof of Theorem~7.2, p.~48]{DasHaconPaun2024}.

These cases cover every irreducible positive-dimensional \(V\).
Additivity over top-dimensional irreducible components gives strict
positivity for every positive-dimensional compact reduced subspace.
We have proved analytic nefness of \(\beta\), the dominating current
\eqref{eq:outer-current-big}, and all positive top intersections.
The normal compact \(Z\) and the smooth real closed \(\beta_0\)
therefore satisfy the three hypotheses of
\cite[Theorem~2.29]{DasHaconPaun2024}.  That theorem gives a positive
smooth representative of \(\beta\).  Its local smooth potentials make
this a K\"ahler form \(\omega_{Z_i}\).  Equation
\eqref{eq:outer-descended-class} gives
\eqref{eq:outer-target-kahler}, and \(Z_i\) is compact K\"ahler.
